%% =================================================================
%% Chen et al. Measurement 264 (2026) 120300.
%% Reconstruction of page 11 (printed p. 11).
%% Contents: Fig. 15 (FPCB defect detection, 8 sub-panels incl.
%% ultrasonic C-scan comparison), Section 4.5 prose covering the
%% effect of material properties and material types.
%% =================================================================

\documentclass[11pt]{article}

\usepackage[margin=0.85in]{geometry}
\usepackage{amsmath, amssymb}
\usepackage{mathrsfs}
\usepackage{xcolor}
\usepackage{fancyhdr}

\pagestyle{fancy}
\fancyhf{}
\lhead{\small Z.~Chen et al.}
\rhead{\small Measurement 264 (2026) 120300}
\cfoot{\small 11 $\to$ \arabic{page}}
\renewcommand{\headrulewidth}{0.4pt}

\setcounter{page}{1}
\setcounter{equation}{1}
\setlength{\parskip}{6pt}
\setlength{\parindent}{0pt}

\begin{document}

% ---------- Fig. 15 placeholder ----------
\begin{center}
\fbox{\begin{minipage}[c][9cm][c]{0.95\textwidth}
\centering
\textbf{Fig. 15.}\ Nondestructive testing of encapsulated flexible
printed circuit board.\\[4pt]
\textcolor{gray}{\footnotesize
\textbf{(a)} A flexible printed circuit board (photograph of the
bare FPCB with a serpentine copper trace pattern, 15 mm scale).\\[3pt]
\textbf{(b)} Flexible printed circuit board encapsulated with a
soft surface layer (photograph of the black encapsulated block).\\[3pt]
\textbf{(c)} Scanning area of the encapsulated FPCB: 28 mm $\times$
10 mm region with a 7.5 mm inner scale bar, showing the trace
pattern through the encapsulation.\\[3pt]
\textbf{(d)} Same as (c) with three injected defects: gap,
scratch, cut (marked by pink dashed boxes).\\[3pt]
\textbf{(e)} 3D imaging of compression-depth distribution when the
tactile tomography system scanned the encapsulated FPCB (c): clean
recovery of the serpentine trace pattern as a rainbow-colored
embossed surface (Z 0--0.5 mm).\\[3pt]
\textbf{(f)} Same as (e) for the FPCB with three defects (d): the
scan shows trace continuity interruptions at the pink-boxed defect
sites.\\[3pt]
\textbf{(g)} Ultrasonic C-scan image of the encapsulated FPCB (c):
blurry greenish map with serpentine outline barely resolved.\\[3pt]
\textbf{(h)} Ultrasonic C-scan image of the FPCB with three defects
(d): similarly low-contrast map with the defect sites indicated but
less distinguishable than in the tactile-tomography scan.]}
\end{minipage}}
\end{center}

\vspace{0.4em}

\textbf{Fig. 15.}\ Nondestructive testing of encapsulated flexible
printed circuit board. (a) A flexible printed circuit board.
(b) A flexible printed circuit board encapsulated with a soft
surface layer. (c) Scanning area of the encapsulated flexible
printed circuit board. (d) Scanning area of the encapsulated
flexible printed circuit board with three defects. (e) 3D imaging
of compression depth distribution when the tactile tomography
system scanned the encapsulated flexible printed circuit board.
(f) 3D imaging of compression depth distribution when the tactile
tomography system scanned the encapsulated flexible printed circuit
board with three defects. (g) Ultrasonic C-scan image of the
encapsulated flexible printed circuit board. (h) Ultrasonic C-scan
image of the encapsulated flexible printed circuit board with three
defects.

\vspace{0.8em}

% ---------- Section 4.5 prose ----------
\subsection*{4.5.\ Effect of materials properties and materials types on the tactile tomography system}

The property of the surface layer materials, especially hardness,
can have an impact on the imaging quality of the tactile tomography
system. As shown in Fig.~13a, the soft surface layer of AB glue was
prepared with different Shore hardness values of 5, 10, 15, and 30
HC, respectively. When the samples were scanned by the tactile
tomography system with a scanning speed and scanning step of 2.4
mm/s \& 0.5 mm under a constant pressure of 4 MPa (Fig.~13b, c, e,
and g), the area where the topographical profile of the spherical
cap model can be reconstructed gradually decreased from the bottom
to the top as Shore hardness of the soft surface layer increased.
Even only the top part of the spherical cap model can be
reconstructed when the Shore hardness of the soft surface layer is
30 HC (Fig.~13g). The reason is that the increase in hardness of
the surface layer will reduce the compression depth of the probe
tip when the samples are scanned by the tactile tomography system
under the same constant pressure, resulting in the inability of
this system to detect and reconstruct the complete internal
structure of the samples. Significantly, the topographical profile
of the spherical cap model can be completely reconstructed and
maintain excellent imaging quality when the probe pressure is
increased for the samples with harder AB glue layer ($P = 6$ MPa
for 10 HC (Fig.~13d), $P = 8$ MPa for 15 HC (Fig.~13f), and $P =
13$ MPa for 30 HC (Fig.~13h)). Although the increase of hardness
of the surface layer, the tactile tomography system can still
maintain a good imaging quality by increasing the probe pressure.
The tactile tomography system can adapt to the change in the
mechanical properties of AB glue with different Shore hardness by
adjusting the probe pressure. Therefore, the tactile tomography
system demonstrates excellent adaptability to the surface layer
materials with different Shore hardness.

In addition to the AB glue, the other two commonly used soft
materials, PDMS and Ecoflex, were also used as the surface layer
to further investigate the versatility of the tactile tomography
system in characterizing samples with different materials
composition. As shown in Fig.~14a, b, and c, the spherical cap
model filled with a soft surface layer of AB glue, PDMS, and
Ecoflex, respectively. These samples were

\end{document}
